% Stacks Project — Tag 09XI
% Lemma 15.11.6 — "Lifting idempotents along henselian pair"
%   (more accurately: characterizations of henselian pairs, including the
%   bijection-on-idempotents criterion that is what 09XI is normally cited
%   for)
% Chapter: More on Algebra (more-algebra.tex)
% Section 15.11 (Henselian pairs)
% Source: https://stacks.math.columbia.edu/tag/09XI
% Source TeX: https://raw.githubusercontent.com/stacks/stacks-project/master/more-algebra.tex
% Fetched: 2026-06-02
% Verbatim excerpt — DO NOT EDIT.
%
% This file vendors the main lemma (09XI) plus the helper lemmas its
% proof directly cites and that are themselves in more-algebra.tex
% (the Stacks-internal cross-references):
%   - lemma-lift-idempotent-upstairs        (more-algebra L1977)
%   - lemma-idempotents-determined-modulo-radical (L2251)
%   - lemma-check-isomorphism-zariski       (L2264)
%   - lemma-helper-finite                   (L2289)
%   - lemma-helper-finite-type              (L2501)
%   - definition-zariski-pair               (L2245)
%   - definition-henselian-pair             (L2393)
%   - lemma-characterize-henselian-pair = TAG 09XI  (L2552)
%
% \ref{...} commands point at labels in other Stacks chapters
% (algebra.tex etc.); to look those up, use the URL pattern
% https://stacks.math.columbia.edu/tag/XXXX with the corresponding
% tag (e.g. lemma-disjoint-decomposition is tag 00EE in algebra.tex).

% ----------------------------------------------------------------------
% From more-algebra.tex, Section 15.10 "Lifting", lines 1977-2017
% ----------------------------------------------------------------------

\begin{lemma}
\label{lemma-lift-idempotent-upstairs}
Let $A$ be a ring, let $I \subset A$ be an ideal.
Let $A \to B$ be an integral ring map.
Let $\overline{e} \in B/IB$ be an idempotent.
Then there exists an \'etale ring map $A \to A'$
which induces an isomorphism $A/I \to A'/IA'$ and an idempotent
$e' \in B \otimes_A A'$ lifting $\overline{e}$.
\end{lemma}

\begin{proof}
Choose an element $y \in B$ lifting $\overline{e}$.
Choose $f \in A[x]$ as in Lemma \ref{lemma-helper-integral} for $y$.
By Lemma \ref{lemma-lift-factorization-easy}
we can find an \'etale ring map $A \to A'$ which induces
an isomorphism $A/I \to A'/IA'$ and such that $f = gh$
in $A[x]$ with $g(x) = x^d \bmod IA'$ and $h(x) = (x - 1)^d \bmod IA'$.
After replacing $A$ by $A'$ we may assume that the factorization
is defined over $A$. In that case we see that
$b_1 = g(y) \in B$ is a lift of $\overline{e}^d = \overline{e}$ and
$b_2 = h(y) \in B$ is a lift of
$(\overline{e} - 1)^d = (-1)^d (1 - \overline{e})^d = (-1)^d(1 - \overline{e})$
and moreover $b_1b_2 = 0$. Thus $(b_1, b_2)B/IB = B/IB$ and
$V(b_1, b_2) \subset \Spec(B)$ is disjoint from $V(IB)$. Since
$\Spec(B) \to \Spec(A)$ is closed (see
Algebra, Lemmas \ref{algebra-lemma-integral-going-up} and
\ref{algebra-lemma-going-up-closed})
we can find an $a \in A$ which maps to an invertible element
of $A/I$ whose image in $B$ lies in $(b_1, b_2)$, see
Lemma \ref{lemma-separate-image-closed-from-closed}.
After replacing $A$ by the localization $A_a$ we get that
$(b_1, b_2) = B$. Then $\Spec(B) = D(b_1) \amalg D(b_2)$;
disjoint union because $b_1b_2 = 0$ and covers
$\Spec(B)$ because $(b_1, b_2) = B$. Let $e \in B$ be the idempotent
corresponding to the open and closed subset $D(b_1)$, see
Algebra, Lemma \ref{algebra-lemma-disjoint-decomposition}.
Since $b_1$ is a lift of $\overline{e}$ and $b_2$ is a
lift of $\pm (1 - \overline{e})$ we conclude that $e$ is
a lift of $\overline{e}$ by the uniqueness statement in
Algebra, Lemma \ref{algebra-lemma-disjoint-decomposition}.
\end{proof}

% ----------------------------------------------------------------------
% From more-algebra.tex, Section 15.11 preamble, lines 2245-2350
% ----------------------------------------------------------------------

\begin{definition}
\label{definition-zariski-pair}
A {\it Zariski pair} is a pair $(A, I)$ such that
$I$ is contained in the Jacobson radical of $A$.
\end{definition}

\begin{lemma}
\label{lemma-idempotents-determined-modulo-radical}
Let $(A, I)$ be a Zariski pair. Then the map from
idempotents of $A$ to idempotents of $A/I$ is injective.
\end{lemma}

\begin{proof}
An idempotent of a local ring is either $0$ or $1$.
Thus an idempotent is determined by the set of maximal ideals
where it vanishes, by
Algebra, Lemma \ref{algebra-lemma-characterize-zero-local}.
\end{proof}

\begin{lemma}
\label{lemma-check-isomorphism-zariski}
Let $(A, I)$ be a Zariski pair. Let $A \to B$ be a flat,
integral, finitely presented ring map such that $A/I \to B/IB$
is an isomorphism. Then $A \to B$ is an isomorphism.
\end{lemma}

\begin{proof}
The ring map $A \to B$ is finite by
Algebra, Lemma \ref{algebra-lemma-characterize-finite-in-terms-of-integral}.
Hence $B$ is finitely presented as an $A$-module by
Algebra, Lemma \ref{algebra-lemma-finite-finitely-presented-extension}.
Hence $B$ is a finite locally free $A$-module by
Algebra, Lemma \ref{algebra-lemma-finite-projective}.
Since the module $B$ has rank $1$
along $V(I)$ (see rank function described in
Algebra, Lemma \ref{algebra-lemma-finite-projective}),
and as $(A, I)$ is a Zariski pair, we conclude that the
rank is $1$ everywhere.
It follows that $A \to B$ is an isomorphism:
it is a pleasant exercise to show that
a ring map $R \to S$ such that $S$ is a locally free $R$-module
of rank $1$ is an isomorphism (hint: look at local rings).
\end{proof}

\begin{lemma}
\label{lemma-helper-finite}
Let $(A, I)$ be a Zariski pair. Let $A \to B$ be a finite ring map.
Assume
\begin{enumerate}
\item $B/IB = B_1 \times B_2$ is a product of $A/I$-algebras
\item $A/I \to B_1/IB_1$ is surjective,
\item $b \in B$ maps to $(1, 0)$ in the product.
\end{enumerate}
Then there exists a monic $f \in A[x]$ with $f(b) = 0$ and
$f \bmod I = (x - 1)x^d$ for some $d \geq 1$.
\end{lemma}

\begin{proof}
By Lemma \ref{lemma-lift-idempotent-upstairs} we can find an
\'etale ring map $A \to A'$ inducing an isomorphism $A/I \to A'/IA'$
such that $B' = B \otimes_A A'$ contains
an idempotent $e'$ lifting the image of $b$ in $B'/IB'$.
Consider the corresponding $A'$-algebra decomposition
$$
B' = B'_1 \times B'_2
$$
which is compatible with the one given in the lemma upon
reduction modulo $I$.
The map $A' \to B'_1$ is surjective modulo $IA'$. By Nakayama's lemma
(Algebra, Lemma \ref{algebra-lemma-NAK})
we can find $i \in IA'$ such that after replacing $A'$
by $A'_{1 + i}$ the map $A' \to B'_1$ is surjective.
Observe that the image $b'_1 \in B'_1$ of $b$
satisfies $b'_1 - 1 \in IB'_1$.
Thus we may pick $a' \in IA'$ mapping to $b'_1 - 1$.
On the other hand, the image $b'_2 \in B'_2$ of $b$ is in $IB'_2$. By
Algebra, Lemma \ref{algebra-lemma-integral-integral-over-ideal}
there exist a monic polynomial
$g(x) = x^d + \sum a'_j x^j$ of degree $d$ with $a'_j \in IA'$ such that
$g(b'_2) = 0$ in $B'_2$. Thus the image $b' = (b'_1, b'_2) \in B'$
of $b$ is a root of the polynomial $(x - 1 - a')g(x)$. We conclude that
$$
(b' - 1)(b')^d \in \sum\nolimits_{j = 0, \ldots, d} IA' \cdot (b')^j
$$
We claim that this implies
$$
(b - 1)b^d \in \sum\nolimits_{j = 0, \ldots, d} I \cdot b^j
$$
in $B$. For this it is enough to see that the ring map
$A \to A'$ is faithfully flat, because the condition is that
the image of $(b - 1)b^d$ is zero in $B/\sum_{j = 0, \ldots, d} Ib^j$
(use Algebra, Lemma \ref{algebra-lemma-faithfully-flat-universally-injective}).
The map $A \to A'$ flat because it is \'etale
(Algebra, Lemma \ref{algebra-lemma-etale}).
On the other hand, the induced map on spectra is open
(see Algebra, Proposition \ref{algebra-proposition-fppf-open} and use
previous lemma referenced) and the image
contains $V(I)$. Since $I$ is contained in the Jacobson radical of $A$
we conclude.
\end{proof}

% ----------------------------------------------------------------------
% From more-algebra.tex, Section 15.11 "Henselian pairs",
% lines 2393-2410 (definition) and 2501-2549 (helper-finite-type)
% ----------------------------------------------------------------------

\begin{definition}
\label{definition-henselian-pair}
A {\it henselian pair} is a pair $(A, I)$ satisfying
\begin{enumerate}
\item $I$ is contained in the Jacobson radical of $A$, and
\item for any monic polynomial $f \in A[T]$ and factorization
$\overline{f} = g_0h_0$ with $g_0, h_0 \in A/I[T]$ monic
generating the unit ideal in $A/I[T]$, there
exists a factorization $f = gh$ in $A[T]$ with $g, h$ monic
and $g_0 = \overline{g}$ and $h_0 = \overline{h}$.
\end{enumerate}
\end{definition}

\begin{lemma}
\label{lemma-helper-finite-type}
Let $(A, I)$ be a pair. Let $A \to B$ be a finite type ring map
such that $B/IB = C_1 \times C_2$ with $A/I \to C_1$ finite.
Let $B'$ be the integral closure of $A$ in $B$.
Then we can write $B'/IB' = C_1 \times C'_2$ such that
the map $B'/IB' \to B/IB$ preserves product decompositions
and there exists a $g \in B'$ mapping to $(1, 0)$ in
$C_1 \times C'_2$ with $B'_g \to B_g$ an isomorphism.
\end{lemma}

\begin{proof}
Observe that $A \to B$ is quasi-finite at every prime of the
closed subset $T = \Spec(C_1) \subset \Spec(B)$ (this follows
by looking at fibre rings, see
Algebra, Definition \ref{algebra-definition-quasi-finite}).
Consider the diagram of topological spaces
$$
\xymatrix{
\Spec(B) \ar[rr]_\phi \ar[rd]_\psi & & \Spec(B') \ar[ld]^{\psi'} \\
& \Spec(A)
}
$$
By Algebra, Theorem \ref{algebra-theorem-main-theorem}
for every $\mathfrak p \in T$ there is a $h_\mathfrak p \in B'$,
$h_\mathfrak p \not \in \mathfrak p$ such that $B'_h \to B_h$ is
an isomorphism. The union $U = \bigcup D(h_\mathfrak p)$ gives an open
$U \subset \Spec(B')$ such that $\phi^{-1}(U) \to U$ is a homeomorphism
and $T \subset \phi^{-1}(U)$. Since $T$ is open in $\psi^{-1}(V(I))$
we conclude that $\phi(T)$ is open in $U \cap (\psi')^{-1}(V(I))$.
Thus $\phi(T)$ is open in $(\psi')^{-1}(V(I))$.
On the other hand, since $C_1$ is finite over $A/I$ it is
finite over $B'$. Hence $\phi(T)$ is a closed subset of $\Spec(B')$
by Algebra, Lemmas \ref{algebra-lemma-going-up-closed} and
\ref{algebra-lemma-integral-going-up}. We conclude that
$\Spec(B'/IB') \supset \phi(T)$ is open and closed. By
Algebra, Lemma \ref{algebra-lemma-disjoint-implies-product}
we get a corresponding product decomposition $B'/IB' = C'_1 \times C'_2$.
The map $B'/IB' \to B/IB$ maps $C'_1$ into $C_1$ and $C'_2$ into $C_2$
as one sees by looking at what happens on spectra (hint: the inverse
image of $\phi(T)$ is exactly $T$; some details omitted).
Pick a $g \in B'$ mapping to $(1, 0)$ in $C'_1 \times C'_2$
such that $D(g) \subset U$; this is possible because $\Spec(C'_1)$
and $\Spec(C'_2)$ are disjoint and closed in $\Spec(B')$ and
$\Spec(C'_1)$ is contained in $U$. Then $B'_g \to B_g$ defines a homeomorphism
on spectra and an isomorphism on local rings (by our choice of $U$ above).
Hence it is an isomorphism, as follows for example from
Algebra, Lemma \ref{algebra-lemma-characterize-zero-local}.
Finally, it follows that $C'_1 = C_1$ and the proof is complete.
\end{proof}

% ======================================================================
% MAIN RESULT: Tag 09XI
% ======================================================================
% From more-algebra.tex, Section 15.11 "Henselian pairs",
% lines 2552-2734.
% ======================================================================

\begin{lemma}
\label{lemma-characterize-henselian-pair}
\begin{reference}
\cite[Chapter XI]{Henselian} and \cite[Proposition 1]{Gabber-henselian}
\end{reference}
Let $(A, I)$ be a pair. The following are equivalent
\begin{enumerate}
\item $(A, I)$ is a henselian pair,
\item given an \'etale ring map $A \to A'$ and an $A$-algebra map
$\sigma : A' \to A/I$, there exists an $A$-algebra map $A' \to A$
lifting $\sigma$,
\item for any finite $A$-algebra $B$ the map $B \to B/IB$ induces
a bijection on idempotents,
\item for any integral $A$-algebra $B$ the map $B \to B/IB$ induces
a bijection on idempotents, and
\item (Gabber) $I$ is contained in the Jacobson radical of $A$ and
every monic polynomial $f(T) \in A[T]$ of the form
$$
f(T) = T^n(T - 1) + a_n T^n + \ldots + a_1 T + a_0
$$
with $a_n, \ldots, a_0 \in I$ and $n \ge 1$ has a root $\alpha \in 1 + I$.
\end{enumerate}
Moreover, in part (5) the root is unique.
\end{lemma}

\begin{proof}
Assume (2) holds. Then $I$ is contained in the Jacobson radical of $A$, since
otherwise there would be a nonunit $f \in A$ congruent to $1$ modulo $I$
and the map $A \to A_f$ would contradict (2). Hence $IB \subset B$
is contained in the Jacobson radical of $B$ for $B$ integral over $A$
because $\Spec(B) \to \Spec(A)$ is closed by
Algebra, Lemmas \ref{algebra-lemma-going-up-closed} and
\ref{algebra-lemma-integral-going-up}.
Thus the map from idempotents of $B$ to idempotents of $B/IB$
is injective by Lemma \ref{lemma-idempotents-determined-modulo-radical}.
On the other hand, since (2) holds, every idempotent
of $B/IB$ lifts to an idempotent of $B$
by Lemma \ref{lemma-lift-idempotent-upstairs}.
In this way we see that (2) implies (4).

\medskip\noindent
The implication (4) $\Rightarrow$ (3) is trivial.

\medskip\noindent
Assume (3). Let $\mathfrak m$ be a maximal ideal and consider the
finite map $A \to B = A/(I \cap \mathfrak m)$. The condition that
$B \to B/IB$ induces a bijection on idempotents implies that
$I \subset \mathfrak m$ (if not, then $B = A/I \times A/\mathfrak m$
and $B/IB = A/I$). Thus we see that $I$ is contained in the Jacobson
radical of $A$. Let $f \in A[T]$ be monic and suppose given a
factorization $\overline{f} = g_0h_0$ with $g_0, h_0 \in A/I[T]$ monic
generating the unit ideal in $A/I[T]$.
Set $B = A[T]/(f)$. Let $\overline{e}$ be the idempotent
of $B/IB$ corresponding to the decomposition
$$
B/IB = A/I[T]/(g_0) \times A/I[T]/(h_0)
$$
of $A$-algebras. Let $e \in B$ be an idempotent lifting $\overline{e}$
which exists as we assumed (3). This gives a product decomposition
$$
B = eB \times (1 - e)B
$$
Note that $B$ is free of rank $\deg(f)$ as an $A$-module.
Hence $eB$ and $(1 - e)B$ are finite locally free $A$-modules.
However, since $eB$ and $(1 - e)B$ have constant rank
$\deg(g_0)$ and $\deg(h_0)$ over $A/I$ we find that the same
is true over $\Spec(A)$. We conclude that
\begin{align*}
f & = \text{CharPol}_A(T : B \to B) \\
& =
\text{CharPol}_A(T : eB \to eB)
\text{CharPol}_A(T : (1 - e)B \to (1 - e)B)
\end{align*}
is a factorization into monic polynomials reducing to the given
factorization modulo $I$. Here $\text{CharPol}_A$ denotes the characteristic
polynomial of an endomorphism of a finite locally free module over $A$.
If the module is free the $\text{CharPol}_A$ is defined as the
characteristic polynomial of the corresponding matrix and in
general one uses Algebra, Lemma \ref{algebra-lemma-standard-covering} to
glue. Details omitted. Thus (3) implies (1).

\medskip\noindent
Assume (1). Let $f$ be as in (5). The factorization of $f \bmod I$
as $T^n$ times $T - 1$ lifts to a factorization $f = gh$ with $g$
and $h$ monic by Definition \ref{definition-henselian-pair}.
Then $h$ has to have degree $1$
and we see that $f$ has a root reducing to $1$ modulo $I$.
Finally, $I$ is contained in the Jacobson radical by
the definition of a henselian pair.
Thus (1) implies (5).

\medskip\noindent
Before we give the proof of the last step, let us show that
the root $\alpha$ in (5), if it exists, is unique. Namely,
due to the explicit shape of $f(T)$, we have
$f'(\alpha) \in 1 + I$ where $f'$ is the derivative of $f$ with
respect to $T$. An elementary argument shows that
$$
f(T) = f(\alpha + T - \alpha) =
f(\alpha) + f'(\alpha) \cdot (T - \alpha) \bmod
(T - \alpha)^2 A[T]
$$
This shows that any other root $\alpha' \in 1 + I$ of $f(T)$
satisfies $0 = f(\alpha') - f(\alpha) = (\alpha' - \alpha)(1 + i)$
for some $i \in I$, so that, since $1 + i$ is a unit in $A$,
we have $\alpha = \alpha'$.

\medskip\noindent
Assume (5). We will show that (2) holds, in other words, that
for every \'etale map $A \to A'$, every section $\sigma : A' \to A/I$
modulo $I$ lifts to a section $A' \to A$.
Since $A \to A'$ is \'etale, the section $\sigma$
determines a decomposition
\begin{equation}
\label{equation-GCHP}
A'/IA' \cong A/I \times C
\end{equation}
of $A/I$-algebras. Namely, the surjective ring map
$A'/IA' \to A/I$ is \'etale by
Algebra, Lemma \ref{algebra-lemma-map-between-etale}
and then we get the desired idempotent by
Algebra, Lemma \ref{algebra-lemma-surjective-flat-finitely-presented}.
We will show that this decomposition lifts to a decomposition
\begin{equation}
\label{equation-GCHP-want}
A' \cong A'_1 \times A'_2
\end{equation}
of $A$-algebras with $A'_1$ integral over $A$. Then $A \to A'_1$
is integral and \'etale and $A/I \to A'_1/IA'_1$ is an isomorphism,
thus $A \to A'_1$ is an isomorphism by
Lemma \ref{lemma-check-isomorphism-zariski}
(here we also use that an \'etale ring map is flat
and of finite presentation, see Algebra, Lemma \ref{algebra-lemma-etale}).

\medskip\noindent
Let $B'$ be the integral closure of $A$ in $A'$. By
Lemma \ref{lemma-helper-finite-type}
we may decompose
\begin{equation}
\label{equation-dec-mod-I}
B'/IB' \cong A/I \times C'
\end{equation}
as $A/I$-algebras compatibly with (\ref{equation-GCHP})
and we may find $b \in B'$ that lifts $(1, 0)$ such that
$B'_b \to A'_b$ is an isomorphism. If the decomposition
(\ref{equation-dec-mod-I}) lifts to a decomposition
\begin{equation}
\label{equation-want-2}
B' \cong B'_1 \times B'_2
\end{equation}
of $A$-algebras, then the induced decomposition
$A' = A'_1 \times A'_2$ will give the
desired (\ref{equation-GCHP-want}): indeed, since $b$ is a unit in $B'_1$
(details omitted),
we will have $B'_1 \cong A'_1$, so that $A'_1$ will be integral over $A$.

\medskip\noindent
Choose a finite $A$-subalgebra $B'' \subset B'$ containing $b$
(observe that any finitely generated $A$-subalgebra of $B'$
is finite over $A$). After enlarging $B''$ we may assume
$b$ maps to an idempotent in $B''/IB''$ producing
\begin{equation}
\label{equation-again-dec-mod-I}
B''/IB'' \cong C''_1 \times C''_2
\end{equation}
Since $B'_b \cong A'_b$ we see that $B'_b$ is of finite type over $A$.
Say $B'_b$ is generated by $b_1/b^n, \ldots, b_t/b^n$ over $A$
and enlarge $B''$ so that $b_1, \ldots, b_t \in B''$. Then
$B''_b \to B'_b$ is surjective as well as injective, hence an isomorphism.
In particular, we see that $C''_1 = A/I$! Therefore $A/I \to C''_1$
is an isomorphism, in particular surjective.
By Lemma \ref{lemma-helper-finite} we can find
an $f(T) \in A[T]$ of the form
$$
f(T) = T^n(T - 1) + a_n T^n + \ldots + a_1 T + a_0
$$
with $a_n, \ldots, a_0 \in I$ and $n \ge 1$ such that $f(b) = 0$.
In particular, we find that $B'$ is a $A[T]/(f)$-algebra.
By (5) we deduce there is a root $a \in 1 + I$ of $f$.
This produces a product decomposition $A[T]/(f) = A[T]/(T - a) \times D$
compatible with the splitting (\ref{equation-dec-mod-I}) of $B'/IB'$.
The induced splitting of $B'$ is then a desired (\ref{equation-want-2}).
\end{proof}
